\chapter{Conception et Architecture du Système}
\label{chap:conception-architecture}

\section*{Introduction du chapitre}
\addcontentsline{toc}{section}{Introduction du chapitre}

Ce chapitre présente l'architecture technique de la solution proposée ainsi que les choix technologiques adoptés. Il décrit également la conception globale du système, les différentes composantes et leurs interactions.

L'architecture a été conçue en tenant compte des contraintes de performance, de sécurité et de maintenabilité, tout en favorisant une extensibilité future.


\section{Choix technologiques}
\label{sec:choix-technologiques}

\subsection{Backend principal : Spring Boot}
\label{subsec:backend-spring}

\textbf{Spring Boot} a été choisi comme framework principal pour le backend de l'application.

\textbf{Justification} :

\begin{itemize}
    \item \textbf{Robustesse et maturité} : framework éprouvé avec un écosystème riche et bien documenté;
    
    \item \textbf{Architecture structurée} : permet une organisation claire en couches métier (contrôleurs, services, repositories);
    
    \item \textbf{Sécurité intégrée} : Spring Security facilite la gestion de l'authentification et de l'autorisation;
    
    \item \textbf{Performance} : excellentes performances pour les applications web d'entreprise;
    
    \item \textbf{Intégration aisée} : compatible avec la plupart des services externes et APIs.
\end{itemize}

Le backend gère l'ensemble de la logique métier, la persistance des données et l'orchestration des appels aux microservices.


\subsection{Microservice IA : FastAPI + Python}
\label{subsec:microservice-fastapi}

Un microservice dédié, développé avec \textbf{FastAPI}, gère l'extraction et le traitement des données basés sur l'intelligence artificielle.

\textbf{Justification} :

\begin{itemize}
    \item \textbf{Rapidité de développement} : FastAPI offre une syntaxe moderne et productive;
    
    \item \textbf{Performance} : une des frameworks les plus rapides du marché Python;
    
    \item \textbf{Écosystème IA} : Python dispose des meilleures bibliothèques pour le machine learning et le traitement NLP;
    
    \item \textbf{Séparation des responsabilités} : isoler la logique IA permet une meilleure scalabilité;
    
    \item \textbf{Documentation automatique} : FastAPI génère automatiquement la documentation \acrfull{api} (Swagger).
\end{itemize}

Ce service expose des endpoints pour l'extraction de CV, l'analyse de texte et la génération de contenu.


\subsection{Base de données : PostgreSQL}
\label{subsec:base-donnees-postgresql}

\textbf{PostgreSQL} a été sélectionné comme système de gestion de base de données principal.

\textbf{Justification} :

\begin{itemize}
    \item \textbf{Fiabilité et stabilité} : SGBDR de production reconnu pour sa robustesse;
    
    \item \textbf{Performance} : excellentes performances avec des données structurées et relationnelles;
    
    \item \textbf{Adéquation au domaine} : le modèle relationnel s'adapte parfaitement au domaine RH;
    
    \item \textbf{Extensibilité} : support des types JSON pour plus de flexibilité;
    
    \item \textbf{Open source et gratuit} : réduction des coûts d'infrastructure.
\end{itemize}

La base de données stocke les profils, les candidatures, les offres d'emploi et la plupart des données applicatives.


\subsection{Gestion des identités : Keycloak}
\label{subsec:keycloak-identities}

\textbf{Keycloak} gère la centralisation des identités et des accès.

\textbf{Justification} :

\begin{itemize}
    \item \textbf{Authentification robuste} : support de multiples protocoles (OAuth2, OpenID Connect, SAML);
    
    \item \textbf{Gestion des rôles} : système flexible de rôles et de permissions;
    
    \item \textbf{Sécurité centralisée} : unique source de vérité pour les identités;
    
    \item \textbf{Standard industriel} : solution respectant les standards de sécurité.
\end{itemize}

Keycloak s'exécute dans un conteneur séparé avec sa propre base de données, garantissant une isolation maximale.


\subsection{Frontend : React}
\label{subsec:frontend-react}

Le frontend utilise \textbf{React}, une bibliothèque JavaScript pour la construction d'interfaces utilisateur dynamiques.

\textbf{Justification} :

\begin{itemize}
    \item \textbf{Dynamique et réactif} : approche basée sur les composants pour une interface fluide;
    
    \item \textbf{Expérience utilisateur} : rendu côté client pour une meilleure responsivité;
    
    \item \textbf{Écosystème riche} : de nombreuses bibliothèques et outils disponibles;
    
    \item \textbf{Maintenabilité} : code structuré et facile à maintenir.
\end{itemize}

React consomme les \acrfull{api} fournies par le backend Spring Boot et affiche les données de manière dynamique.


\subsection{Orchestration des workflows : n8n}
\label{subsec:orchestration-n8n}

\textbf{n8n} est une plateforme d'orchestration de workflows open source permettant d'automatiser les processus complexes.

\textbf{Justification -- Point fort du projet} :

\begin{itemize}
    \item \textbf{Automatisation visuelle} : création de workflows sans code (low-code);
    
    \item \textbf{Connectivité étendue} : intégration avec des centaines de services et \acrfull{api};
    
    \item \textbf{Logique intelligente} : branchements conditionnels, boucles, transformations de données;
    
    \item \textbf{Orchestration métier} : permet de lier les différentes étapes du processus RH;
    
    \item \textbf{Scalabilité} : ajout facile de nouveaux processus sans modification du code.
\end{itemize}

n8n constitue le cœur de l'automatisation du système, orchestrant les interactions entre le backend, le microservice IA et les systèmes externes.


\subsection{Intelligence artificielle : Groq (LLaMA 3.3)}
\label{subsec:ia-groq}

Le système utilise l'API \textbf{Groq} pour accéder au modèle \textbf{LLaMA 3.3}, un grand modèle de langage haute performance.

\textbf{Justification} :

\begin{itemize}
    \item \textbf{Performance} : Groq offre une inférence très rapide (40\%+ plus rapide que les alternatives);
    
    \item \textbf{Qualité} : LLaMA 3.3 offre des résultats de haute qualité pour l'extraction et l'analyse de texte;
    
    \item \textbf{Coût} : tarification compétitive avec une facturation à l'utilisation;
    
    \item \textbf{Fiabilité} : API bien documentée et stable.
\end{itemize}

\textbf{Note} : \textit{Une alternative envisagée est Ollama, permettant de déployer les modèles en local pour plus de sécurité et de confidentialité des données.}


\subsection{Conteneurisation : Docker}
\label{subsec:containerisation-docker}

L'ensemble de la plateforme est containerisée avec \textbf{Docker}.

\textbf{Justification} :

\begin{itemize}
    \item \textbf{Isolation} : chaque service s'exécute dans son propre conteneur avec son environnement;
    
    \item \textbf{Déploiement facile} : garantit la cohérence entre développement et production;
    
    \item \textbf{Scalabilité} : facilite la réplication des services;
    
    \item \textbf{Maintenabilité} : gestion centralisée des versions et des dépendances via \texttt{docker-compose}.
\end{itemize}


\section{Architecture globale du système}
\label{sec:architecture-globale}

\subsection{Vue d'ensemble}
\label{subsec:vue-ensemble}

L'architecture générale du système suit un pattern microservices avec orchestration de workflows. La figure \ref{fig:architecture-overview} présente les principaux composants et leurs interactions.

\begin{figure}[!ht]
    \centering
    \fbox{\parbox{\textwidth}{
        \textit{Remarque : Un diagramme d'architecture détaillé devrait être inséré ici, montrant :}
        \begin{itemize}
            \item Frontend React (port 3000)
            \item Backend Spring Boot (port 8080)
            \item Microservice FastAPI (port 8000)
            \item n8n (port 5678)
            \item PostgreSQL (conteneur)
            \item Keycloak (conteneur)
            \item Les connexions API entre ces composants
        \end{itemize}
    }}
    \caption{Vue d'ensemble de l'architecture du système}
    \label{fig:architecture-overview}
\end{figure}

Les composants principaux sont :

\begin{itemize}
    \item \textbf{Frontend React} : interface utilisateur web, consomme les API du backend;
    
    \item \textbf{Backend Spring Boot} : logique métier, gestion des données, coordination;
    
    \item \textbf{Microservice FastAPI} : traitement IA, extraction de CV, analyse;
    
    \item \textbf{n8n} : orchestration des workflows automatisés;
    
    \item \textbf{PostgreSQL} : stockage des données applicatives;
    
    \item \textbf{Keycloak} : gestion des identités et authentification;
    
    \item \textbf{Services externes} : API Groq (IA), Google Calendar, services d'email.
\end{itemize}


\subsection{Architecture logique}
\label{subsec:architecture-logique}

L'architecture logique du système est organisée selon une structure en couches :

\begin{enumerate}
    \item \textbf{Couche de présentation} (Frontend React)
    \begin{itemize}
        \item Affichage des interfaces utilisateur,
        \item Gestion des interactions utilisateur,
        \item Consommation des API backend.
    \end{itemize}
    
    \item \textbf{Couche métier} (Backend Spring Boot)
    \begin{itemize}
        \item Logique applicative,
        \item Validation des données,
        \item Gestion des règles métier,
        \item Coordination entre services.
    \end{itemize}
    
    \item \textbf{Couche de traitement IA} (Microservice FastAPI)
    \begin{itemize}
        \item Extraction d'informations des CV,
        \item Analyse de texte,
        \item Appels aux modèles IA (Groq),
        \item Transformation et normalisation des données.
    \end{itemize}
    
    \item \textbf{Couche d'orchestration} (n8n)
    \begin{itemize}
        \item Workflows automatisés,
        \item Chaînage des processus,
        \item Intégration avec les services externes.
    \end{itemize}
    
    \item \textbf{Couche de données} (PostgreSQL, Keycloak)
    \begin{itemize}
        \item Persistance des données applicatives,
        \item Gestion des identités,
        \item Contrôle d'accès.
    \end{itemize}
\end{enumerate}


\subsection{Architecture physique (déploiement)}
\label{subsec:architecture-physique}

Sur le plan du déploiement, le système utilise une architecture containerisée composée de :

\begin{figure}[!ht]
    \centering
    \fbox{\parbox{\textwidth}{
        \textit{Remarque : Un diagramme de déploiement Docker devrait être inséré ici, montrant}
        \begin{itemize}
            \item Conteneurs Docker pour n8n, PostgreSQL, Keycloak,
            \item Services locaux pour Spring Boot et FastAPI,
            \item Réseaux Docker internes,
            \item Volumes persistants.
        \end{itemize}
    }}
    \caption{Architecture physique du déploiement}
    \label{fig:architecture-physique}
\end{figure}

\textbf{Conteneurs Docker} :

\begin{itemize}
    \item \textbf{n8n container} : orchestration des workflows;
    
    \item \textbf{PostgreSQL container} : base de données persistante;
    
    \item \textbf{Keycloak container} : gestion des identités.
\end{itemize}

\textbf{Services locaux} (développement) :

\begin{itemize}
    \item \textbf{Spring Boot} : backend principal;
    
    \item \textbf{FastAPI} : microservice IA;
    
    \item \textbf{React dev server} : frontend.
\end{itemize}

\textbf{Volumes persistants} :

\begin{itemize}
    \item Data PostgreSQL : sauvegarde des données,
    
    \item Configurations n8n : workflows et configurations.
\end{itemize}


\section{Conception du système}
\label{sec:conception-systeme}

\subsection{Diagramme de cas d'utilisation}
\label{subsec:diagramme-uc-chap2}

Le diagramme de cas d'utilisation présenté au chapitre précédent (section \ref{subsec:diagramme-cas-utilisation}) fournit une base pour la conception détaillée du système. Les cas d'utilisation identifiés ont guidé la structuration des modules et des interfaces.


\subsection{Diagramme de classes}
\label{subsec:diagramme-classes}

Les principales entités du domaine RH sont modélisées par les classes suivantes :

\begin{figure}[!ht]
    \centering
    \fbox{\parbox{\textwidth}{
        \textit{Remarque : Un diagramme de classes UML détaillé devrait être inséré ici, montrant :}
        \begin{itemize}
            \item Classe User (avec héritage Admin, RH, Candidat)
            \item Classe Candidate (profil candidate)
            \item Classe CV (structure du CV)
            \item Classe JobOffer (offres d'emploi)
            \item Classe Application (candidatures)
            \item Classe Template (templates CV)
            \item Classe EvaluationScore (scoring)
            \item Les associations et cardinalités
        \end{itemize}
    }}
    \caption{Diagramme de classes du système}
    \label{fig:classe-diagram}
\end{figure}

Les principales entités sont :

\begin{enumerate}
    \item \textbf{User} : représente un utilisateur du système (admin, recruteur, candidat);
    
    \item \textbf{Candidate} : profil d'un candidat, interne ou externe;
    
    \item \textbf{CV} : document d'un candidat, peut exister en plusieurs versions;
    
    \item \textbf{JobOffer} : offre d'emploi publiée;
    
    \item \textbf{Application} : candidature d'un candidat à une offre;
    
    \item \textbf{Template} : modèle de présentation pour les CV;
    
    \item \textbf{EvaluationScore} : score automatique d'une candidature.
\end{enumerate}


\subsection{Diagrammes de séquence}
\label{subsec:diagrammes-sequence}

Deux scénarios clés illustrent les interactions principales du système :

\subsubsection{Scénario 1 : Import et Extraction de CV}
\label{subsubsec:scenario1-import-cv}

Ce scénario décrit le flux complet d'importation et d'extraction automatique d'un CV :

\begin{figure}[!ht]
    \centering
    \fbox{\parbox{\textwidth}{
        \textit{Remarque : Un diagramme de séquence UML devrait être inséré ici, montrant :}
        \begin{enumerate}
            \item L'utilisateur RH importe un CV,
            \item Le backend reçoit le fichier,
            \item Appel au microservice FastAPI,
            \item FastAPI envoie le document à Groq,
            \item Groq extrait les informations,
            \item FastAPI retourne les données extraites,
            \item Le backend stocke les données en base,
            \item Réponse affichée à l'utilisateur.
        \end{enumerate}
    }}
    \caption{Diagramme de séquence : Import et extraction de CV}
    \label{fig:sequence-import-cv}
\end{figure}

\textbf{Étapes du processus} :

\begin{enumerate}
    \item L'utilisateur RH télécharge un CV via l'interface React,
    \item Le backend Spring Boot reçoit le fichier,
    \item Un appel \acrfull{api} est effectué vers le microservice FastAPI,
    \item FastAPI traite le document et appelle l'API Groq,
    \item Groq retourne les informations extraites au format JSON,
    \item FastAPI valide et normalise les données,
    \item Le backend stocke les données en PostgreSQL,
    \item L'interface affiche le CV importé avec les données extraites.
\end{enumerate}


\subsubsection{Scénario 2 : Candidature, Scoring et Communication}
\label{subsubsec:scenario2-candidature-scoring}

Ce scénario illustre le traitement d'une candidature, du scoring automatique à la communication avec le candidat :

\begin{figure}[!ht]
    \centering
    \fbox{\parbox{\textwidth}{
        \textit{Remarque : Un diagramme de séquence UML devrait être inséré ici, montrant :}
        \begin{enumerate}
            \item Un candidat postule à une offre,
            \item La candidature est créée en base,
            \item Un workflow n8n est déclenché,
            \item Appel au scoring IA,
            \item Résultat du scoring stocké,
            \item Génération d'un email de confirmation,
            \item Planification d'un rendez-vous si score élevé,
            \item Notification au recruteur.
        \end{enumerate}
    }}
    \caption{Diagramme de séquence : Candidature et scoring}
    \label{fig:sequence-candidature}
\end{figure}

\textbf{Étapes du processus} :

\begin{enumerate}
    \item Un candidat remplit et soumet une candidature via le site public,
    \item Le backend crée un enregistrement \texttt{Application} en base de données,
    \item Un workflow n8n est déclenché automatiquement,
    \item Le workflow appelle le microservice FastAPI pour le calcul du score,
    \item Le score est calculé et sauvegardé,
    \item Si le score dépasse un seuil, un email de confirmation est envoyé au candidat,
    \item Une proposition de rendez-vous est automatiquement générée et envoyée,
    \item Le recruteur reçoit une notification de la nouvelle candidature.
\end{enumerate}


\section{Conception des workflows intelligents}
\label{sec:conception-workflows}

\subsection{Rôle de n8n}
\label{subsec:role-n8n}

n8n joue un rôle central dans l'automatisation du système. Ses responsabilités principales sont :

\begin{itemize}
    \item \textbf{Orchestration des processus} : coordination des différentes étapes du workflow RH;
    
    \item \textbf{Automatisation} : exécution de tâches sans intervention manuelle;
    
    \item \textbf{Connexion entre services} : liaison du backend, du microservice IA et des services externes;
    
    \item \textbf{Transformation de données} : adaptation des formats de données entre les services.
\end{itemize}

Les workflows n8n offrent une flexibilité maximale et peuvent être modifiés rapidement sans redéployer le code.


\subsection{Exemples de workflows}
\label{subsec:exemples-workflows}

\subsubsection{Workflow 1 : Import et Extraction de CV}
\label{subsubsec:workflow1-import-extraction}

\begin{figure}[!ht]
    \centering
    \fbox{\parbox{\textwidth}{
        \textit{Remarque : Un diagramme de workflow n8n devrait être inséré ici, montrant :}
        \begin{enumerate}
            \item Trigger : appel API du backend,
            \item Nœud de validation du fichier,
            \item Nœud d'appel à FastAPI,
            \item Nœud de traitement du résultat,
            \item Nœud de mise à jour de la base de données,
            \item Nœud de réponse au backend.
        \end{enumerate}
    }}
    \caption{Workflow n8n : Import et extraction de CV}
    \label{fig:workflow-import-cv}
\end{figure}

Ce workflow automatise l'importation et l'extraction des CV :

\begin{enumerate}
    \item \textbf{Trigger API} : réception d'un nouveau CV,
    \item \textbf{Validation du fichier} : vérification du format et de la taille,
    \item \textbf{Appel à FastAPI} : envoi du fichier pour extraction,
    \item \textbf{Traitement du résultat} : validation et normalisation des données extraites,
    \item \textbf{Mise à jour de la base} : sauvegarde des informations en PostgreSQL,
    \item \textbf{Réponse au backend} : notification de la fin du traitement.
\end{enumerate}


\subsubsection{Workflow 2 : Traitement des Candidatures}
\label{subsubsec:workflow2-candidatures}

\begin{figure}[!ht]
    \centering
    \fbox{\parbox{\textwidth}{
        \textit{Remarque : Un diagramme de workflow n8n devrait être inséré ici, montrant :}
        \begin{enumerate}
            \item Trigger : nouvelle candidature,
            \item Nœud de récupération des données du candidat,
            \item Nœud d'appel au scoring IA,
            \item Nœud de branchement conditionnel (score élevé/faible),
            \item Nœud d'envoi d'email,
            \item Nœud de planification du rendez-vous,
            \item Nœud de notification au recruteur.
        \end{enumerate}
    }}
    \caption{Workflow n8n : Traitement des candidatures}
    \label{fig:workflow-candidatures}
\end{figure}

Ce workflow gère automatiquement le traitement des candidatures :

\begin{enumerate}
    \item \textbf{Trigger} : une nouvelle candidature est reçue,
    \item \textbf{Récupération des données} : fetch du profil du candidat et des exigences de l'offre,
    \item \textbf{Calcul du score} : appel au microservice FastAPI pour le scoring IA,
    \item \textbf{Branchement conditionnel} : deux chemins selon le score obtenu,
    \item \textbf{Score élevé} : envoi d'un email personnalisé et proposition de rendez-vous,
    \item \textbf{Score faible} : archivage de la candidature avec notification optionnelle,
    \item \textbf{Notification au recruteur} : mise à jour du tableau de bord.
\end{enumerate}


\subsection{Intégration de l'IA dans les workflows}
\label{subsec:integration-ia-workflows}

L'intelligence artificielle est intégrée dans les workflows de la manière suivante :

\begin{itemize}
    \item \textbf{Appels à Groq} : le microservice FastAPI fait appel à l'API Groq (LLaMA 3.3);
    
    \item \textbf{Validation via Pydantic} : les résultats sont validés à l'aide de modèles Pydantic;
    
    \item \textbf{Traitement JSON} : les données sont transformées au format JSON pour la transmission;
    
    \item \textbf{Mémorisation des résultats} : les résultats de l'IA sont cachés pour améliorer les performances.
\end{itemize}

Cette approche garantit la fiabilité et la cohérence du traitement IA au sein des workflows.


\section{Sécurité du système}
\label{sec:securite-systeme}

\subsection{Authentification}
\label{subsec:authentification}

L'authentification est gérée centralement par \textbf{Keycloak}, qui supporte plusieurs protocoles de sécurité :

\begin{itemize}
    \item \textbf{OAuth2} : pour l'authentification par des tiers (Google, LinkedIn);
    
    \item \textbf{OpenID Connect} : pour l'authentification fédérée;
    
    \item \textbf{Identifiants locaux} : username/password sécurisés avec hachage bcrypt.
\end{itemize}

Tous les utilisateurs doivent s'authentifier avant d'accéder au système. Les tokens JWT sont utilisés pour les appels \acrfull{api} ultérieurs.


\subsection{Autorisation}
\label{subsec:autorisation}

Les rôles et permissions sont définis dans Keycloak :

\begin{itemize}
    \item \textbf{Admin} : accès complet à tous les modules;
    
    \item \textbf{RH/Recruteur} : accès aux candidatures, offres, profils;
    
    \item \textbf{Candidat} : accès limité au profil personnel et aux offres publiques.
\end{itemize}

Le backend vérifie les permissions avant chaque action sensible.


\subsection{Protection des données}
\label{subsec:protection-donnees}

Plusieurs mesures protègent les données sensibles :

\begin{itemize}
    \item \textbf{Connexions HTTPS} : encryption du trafic en transit;
    
    \item \textbf{Base de données chiffrée} : PostgreSQL peut être configuré avec du chiffrement;
    
    \item \textbf{Isolation des services} : réseau Docker fermé, pas d'accès direct;
    
    \item \textbf{Validation des entrées} : prévention des injections SQL et XSS;
    
    \item \textbf{Tokens courts} : expiration rapide des tokens d'authentification.
\end{itemize}


\section{Gestion des données}
\label{sec:gestion-donnees}

\subsection{Modélisation de la base de données}
\label{subsec:modelisation-bd}

La base de données PostgreSQL contient les principales tables suivantes :

\begin{figure}[!ht]
    \centering
    \fbox{\parbox{\textwidth}{
        \textit{Remarque : Un diagramme \acrshort{er} (Entity-Relationship) détaillé devrait être inséré ici, montrant :}
        \begin{itemize}
            \item Tables : users, candidates, cvs, job\_offers, applications, templates, evaluation\_scores,
            \item Clés primaires et étrangères,
            \item Cardinalités des relations,
            \item Types de données des colonnes.
        \end{itemize}
    }}
    \caption{Diagramme de la base de données}
    \label{fig:schema-bd}
\end{figure}

Les tables principales et leurs rôles :

\begin{itemize}
    \item \textbf{users} : comptes utilisateurs et rôles;
    
    \item \textbf{candidates} : profils des candidats (internes et externes);
    
    \item \textbf{cvs} : documents CV et leurs versions;
    
    \item \textbf{job\_offers} : offres d'emploi publiées;
    
    \item \textbf{applications} : candidatures soumises;
    
    \item \textbf{templates} : templates de présentation des CV;
    
    \item \textbf{evaluation\_scores} : résultats du scoring automatique;
    
    \item \textbf{communications} : logs des emails et messages envoyés.
\end{itemize}


\subsection{Flux de données}
\label{subsec:flux-donnees}

Le flux général de circulation des données dans le système est :

\begin{enumerate}
    \item \textbf{Importation} : un CV est téléchargé via le frontend React,
    
    \item \textbf{Traitement IA} : extraction des données via FastAPI et Groq,
    
    \item \textbf{Persistance} : les données sont stockées en PostgreSQL,
    
    \item \textbf{Affichage} : le frontend récupère et affiche les données via les API du backend,
    
    \item \textbf{Automatisation} : n8n utilise les données pour orchestrer des workflows,
    
    \item \textbf{Intégration externe} : envoi de données vers Google Calendar, services d'email, etc.
\end{enumerate}

Cette architecture en couches garantit une séparation claire des responsabilités et une scalabilité maximale.


\section{Conclusion}
\label{sec:conclusion-chap2}

Ce chapitre a présenté l'architecture globale du système et les choix techniques adoptés. L'approche microservices avec orchestration par n8n offre une flexibilité maximale et une scalabilité excellente.

Les technologies choisies (Spring Boot, FastAPI, PostgreSQL, Keycloak, React, n8n) constituent un stack technique moderne et performant, adaptée aux objectifs du projet.

Le chapitre suivant détaillera la mise en œuvre du projet à travers les différents sprints de développement, montrant comment cette architecture a été concrétisée en code et en fonctionnalités opérationnelles.

\newpage
